%% kapa: exam=2 round=37 year=2026 period=2 subject=appraisal_theory subject_ko=감정평가이론 seq=1 number=1 points=40
%% source: Q-Net 감정평가사 기출문제 (cst00309) artlSeq=5268128 file=raw/1008/5268128/26년 제37회 감정평가사 2차 (2교시_감정평가이론).pdf
%% conversion: kordoc 4.24.51 -> tools/gen2.py (auto); review=proofread-gpt-terra; flags=none
\begin{longquestion}{【문제 1】}
다음 자료를 참고하여 물음에 답하시오. (40점) \par
\begin{framedtext}
1. 대상토지 \par
1) 소재지 및 지번: A도 B시 C동 100번지, 지목: 답, 면적: 1,000 m\textsuperscript{2} \par
2) 토지이용계획: 제2종일반주거지역, 도시ㆍ군계획시설(도로: 중로 1류)에 \par
20\% 저촉하며, 도시ㆍ군계획시설(도로) 결정의 고시일은 2010년 7월 7일임 \par
3) 임대차 관계: 토지소유자는 乙과 대상토지에 대해 임대차계약을 체결하고 그 \par
계약에 대해 관할 동장의 확인을 받음 (임대차기간: 2020. 03. 03. ∼ 2030. 03. 02.) \par
4) 이용상황: 乙이 판매 목적의 사과나무 묘목을 재배 중임 \par
5) 구분지상권: 대상토지의 북서측 상공으로 고압(154 kV) 송전선이 통과하며, 2000년 7월 7일 구분지상권이 설정됨 \par
6) 표준지: 대상토지는 표준지로서 2026년 공시지가는 250,000 원/m\textsuperscript{2} 임 \par
2. 주위환경 \par
1) 주위는 남측하향 완경사지대로서 대부분 농지로 이용 중이고, 드문드문 \par
주택이 소재함 \par
2) 북측 인근은 야산지대이며, 남측 인근은 기존 주택지대임 \par
3) 동측 인근에「택지개발촉진법」에 의한 택지개발사업(부지면적: 200,000 m\textsuperscript{2})이 시행 중이며, 사업시행자는 2026년 3월 3일 보상계획 공고를 함 \par
3. 법원은 대상토지에 대해 감정평가사 丙에게 2026년 7월 1일 경매목적의 \par
감정평가를 의뢰함 \par
\end{framedtext}
\medskip\noindent 물음 1) 丙은 대상토지에 대해 지역분석을 하고 있다. 부동산감정평가에 있어 지역분석의 목적과 절차에 근거하여 검토해야 할 사항을 쓰시오. (20점) \par
\medskip\noindent 물음 2) 丙은 대상토지를 비교표준지로 선정하여 공시지가기준법으로 감정평가하고 있다. 선정된 표준지공시지가에 대해 표준지공시지가의 평가기준에 근거하여 검토해야 할 사항을 쓰시오. (20점) \par
\end{longquestion}
